% 原书第 139 页（PDF 第 152 页）
\BookPage{139}{152}
\BookHeading{黄河远上白云间}{黄河远上白云间\\[3pt]—关于河源认识之探讨}{林天人*}
\BookSubheading{一、楔子：院藏黄河图河源的描绘}

国立故宫博物院所藏文物来自不同单位，文物类大部分承接自清宫遗物；但文献舆图类则来自国立北平图书馆，而这部分藏品是故宫重要的收藏项目。国立故宫博物院典藏的古舆图，总计近千件；不过这项统计未必准确，因为典藏地图中常有若干幅合并一件计或一本地图册页中附带多幅，同时也不包括附在奏折中的地图附件。2

故宫所藏地图除了历代古籍上的刻绘外，单幅、单件或长卷的作品，大抵为明、清时期所绘制的居多，其中清代绘制的数量高过明代所绘；但由于清宫入藏地图，编录并不确实，因此很多藏品年代失据。北京第一历史档案馆邹爱莲认为：「清宫各类舆图，虽然清朝时一大部分已编有目录，但大多数没有标着图件形成的具体年代，没有绘制者，加之有的折图分家，和后期征集进馆之图年代考证不确，利用起来

\vfill
\begin{BookNotes}
\BookNote{*}{作者任职于国立故宫博物院。}
\BookNote{1.}{不过，北平图书馆馆藏地图亦来自清宫内阁大库及后来新购；关于北平图书馆所藏地图数量与目录，可参王庸，〈国立北平图书馆清内阁大库舆图目录〉及〈国立北平图书馆特藏新购地图目录〉，载于《国立北平图书馆馆刊》第6卷，第4号，（台北:学生书局重刊，民56)。}
\BookNote{2.}{国立故宫博物院的典藏中，传统舆图类有983件；但这个统计数字若是一幅一幅的计算，未必准确，因为常有若干幅合并一件计或一本地图册页中附带多幅，同时也不包括附在奏折中的地图附件。目前故宫典藏地图中，有以件、幅、卷、页、折、册、帧计；如果以一张地图算一件的话，地图的总数超过数千件。目前故宫舆图类资料已申请国科会数位计划整理中。关于国立故宫博物院舆图入藏方式与所藏地图种类，可参卢雪燕，地理之学非图不明—故宫所藏古地图探微〉，载于《中华民国第七届地图学术研讨会特刊》，（台北，2007）及李孝聪，〈国立故宫博物院图书文献处藏清代舆图的初步整理与认识〉，载于《故宫学术季刊》，第25卷，第1期。卢文另提及故宫对古地图的保存状况；李文则提出藏图重新归类整理的建议，可以参考。}
\end{BookNotes}

% 原书第 140 页（PDF 第 153 页）
\BookPage{140}{153}
\noindent 多有不便。」3邹氏所谓，基本上是中国传统舆图共同的问题。

院藏共有近五十幅以「黄河」为名或与黄河相关的舆图；其中包括卷轴、册页与军机、宫中档折件附图，但不包括古籍中绘制黄河的插图或图中虽牵涉到黄河，但题署无「黄河」之舆图。而这批黄河图中仅少数三、五幅描绘到河源地区；但却没有一幅是以描绘河源地区为主的舆图。推敲原因，或者军机、宫中档折件的黄河图，都是黄河河患与生民攸关的议题；而长幅卷轴的黄河图，亦都是治河相关的描述。因此，院藏与黄河有关的舆图，河源地区大都是附带的部分。但是耐人寻味的是，院藏黄河图中每一幅所绘到的河源均不一致，大概是绘图者对河源地区的认知也不一致。在舆图中，有将星宿海绘成一大片湖泊状；亦有将星宿海绘成藕节状的河源地形。如题署〈黄河发源至河南省各厅工程情形图〉（登录号:平图020882）为名的黄河图册页，将河源放在星宿海，接着黄河河流中注记文字「头巨泽」、「二巨泽」、「三巨泽」等；图中没注其他文字，祇标注地名。

另一幅题署〈六省黄河埽坝情形图〉（登录号:平图020884）为名的黄河图册页，将黄河源置于「巴颜喀拉山」；图中附上注记文字，曰：

\begin{BookQuote}
黄河出今西番巴颜喀拉山东，名阿尔坦河。东北流三千余里，合鄂敦塔拉诸泉源，大小千百泓，错列如星，汇为查灵、鄂灵二湖，各周围三百余里，东西相去五十里。…经积石山至兰州府属河州境，入中国界。
\end{BookQuote}

另外，题为〈黄河本未全图〉（登录号:平图020873）及〈黄河全图〉（登录号:平图020871）都是以「星宿海」为黄河源。其实，星宿海的范围，至今尚未明确；但从元代以后，真正登入河源实地勘察后，以星宿海当河源的说法逐渐被人们所接受，因此星宿海也渐被当作河源的代名词。不过，真正的河源尚在星宿海以西的河源支河。

河源的问题相当复杂，其中原因固然是在历史上很长的一段时期，河源地区为少数民族所控制，中原王朝探源不易；另亦可能受到经籍先入为主的影响。因此，长久以来关于河源错误的讹传，便代代延续下来。

\vfill
\begin{BookNotes}
\BookNote{3.}{参邹爱莲，〈中国第一历史档案馆藏清宫舆图〉，载于《第三届清代档案国际学术研讨会—清代档案整理与馆际合作会议论文》，（台北，2006)，页308。}
\end{BookNotes}

% 原书第 141 页（PDF 第 154 页）
\BookPage{141}{154}
\BookArt{p154-01.jpg}{0.990}{0.343}

\BookArt{p154-02.jpg}{0.990}{0.344}


% 原图说　院藏《黄河发源至河南省各厅工程情形图》（登录号：平图020882）

% 原图说　院藏《六省黄河埽坝情形图》（登录号：平图020884）

% 原书第 142 页（PDF 第 155 页）
\BookPage{142}{155}
\BookSubheading{二、关于河源：一些历史的认识}

黄河对于传统的中国有特殊的意义，它不但被喻为华夏子孙「母亲的河」，同时也被认为是中华民族文化起源的摇篮。但是黄河为患几乎与中国历史相终始；因此，历代君臣百姓对于黄河一直抱持着又敬且畏的态度。

面对这么一条蜿蜒緜长而带有象征意义的河流，历史上留有大量的文献记载与描述各河段；特别是中、下游地区与广大人民的生活息息相关，因此记录河道的各种状况，更是史不绝书。但相对于河源的记载，史料却寥落稀少且多臆断之词；原因是河源固然探勘不易，其次河源地区长期以来被中原视为边陲，而为少数民族所控制有关。历史上最早居住河源地区为古羌族，4他们对于河源的认识并没有留下直接的文字记录。因此，外人对于这个地区的了解，恐怕需借由实际走访探勘得知。杜佑《通典》卷一七四论及「河源」，谓：

\begin{BookQuote}
且汉时群羌种众虽多，不相统一，未为强国，汉家或未尝遣使西南羌中，或未
\end{BookQuote}

知自有河也。大概中原人士在此之前对于河源的认识，都从零星的文献拼凑出来。因此，长久以来对于河源的了解，存在某些积习甚久的臆说。

除了亡佚的《禹本纪》若干文字仍被征引外，现存历史上对于河源较早的记载大概推《山海经》及《尔雅》等先秦文献。《山海经·西山经》云：

\begin{BookQuote}
昆仑之丘，…河水出焉，而东注于无达。
\end{BookQuote}

《尔雅·释水》亦曰：

\begin{BookQuote}
河出昆仑虚，色白。所渠并千七百，一川色黄，百里一小曲，千里一曲一直。
\end{BookQuote}

《山海经》且不止一次提到「昆仑之丘」或「昆仑之墟」，但都无法言明何处；5祇

\vfill
\begin{BookNotes}
\BookNote{4.}{羌戎即《后汉书·西羌传》所云居于「赐支河」滨的西羌；也是《尚书·禹贡》所谓的「析支」。〈西羌传〉云：「西羌之本，出自三苗，姜姓之别也。其国近南岳。及舜流四凶，徙之三危，河关之西南羌地是也。滨于赐支，至乎河首，挠地千里。赐支者，〈禹贡〉所谓析支者也。」秦厉公时(前461)，羌人无弋爰剑「亡入三河间」。李贤注谓：「三河，即黄河、赐支河、湟河也。}
\BookNote{5.}{《山海经》有两昆仑，一在中国西北，帝下之都，河水出其东北者隅是也；一在西南之海，西王母所居，弱水之渊环其下者是也。}
\end{BookNotes}

% 原书第 143 页（PDF 第 156 页）
\BookPage{143}{156}
\noindent 在〈海内西经〉提及「昆仑山在西胡西。」因此，后世祇能据此猜测「昆仑疑当在西域，而不知其在何国」。6

虽然「昆仑」之地不详，但其说影响深远；传世的文献几乎不疑也不辨其说，且在此基础下附会增益，将「昆仑」之地妄加臆度揣测。几乎把「河出昆仑」视为确切不移之论，同时纷纷推敲昆仑位于何处。《水经》引《禹本纪》断昆仑离中土数万里。7《水经·河水注一》说：

\begin{BookQuote}
昆仑墟在西北，…去嵩高五万里，地之中也。
\end{BookQuote}

《水经·河水注二》引《山海经》谓河源出于昆仑，云：

\begin{BookQuote}
…敦薨之水，其水出焉耆之北。敦薨之山，在匈奴之西，乌孙之东。《山海经》曰：敦薨之山，…敦薨之水出焉，而西流注于渤泽。出于昆仑之东北隅，实惟河源者也。8
\end{BookQuote}

《水经》以出于昆仑山之东北的敦薨之水为河源。较早而明确指出昆仑山的位置及河源分布状况为《新唐书·吐蕃传（下）》，曰：

\begin{BookQuote}
穆宗长庆元年（821），以大理卿刘元鼎为吐蕃会盟使，…河之上流，繇洪济梁西南行二千里，水益狭，春可涉，秋夏乃胜舟。其南三百里三山，中高而四下，曰紫山，直大羊同国，古所谓昆仑者也，虏曰闷摩黎山，东距长安五千里，河源其间，流澄缓下，稍合众流，色赤，行益远，它水并注则浊，故举世谓西戎地曰河湟。河源东北直莫贺延碛尾殆五百里，碛广五十里，北自沙州，西南入吐谷浑浸狭，故号碛尾。隐测其地，盖剑南之西。元鼎所经见大略如
\end{BookQuote}

此。

\vfill
\begin{BookNotes}
\BookNote{6.}{参胡渭着、邹逸麟整理，《禹贡锥指》十三上，(上海：古籍出版社，1996版)，页427。胡渭云:「然自汉以前，但知昆仑在中国西北，终未明在外国某地也。」，页426。}
\BookNote{7.}{《水经注》的臆断，杜佑认为出于《禹本纪》；《通典》卷一七四「风俗」曰:「《水经》所云河出昆仑山者，宜出于《禹本纪》。」杨守敬亦说：「此沿《禹本纪》之误，古时中外未通，传闻异辞，转相傅会，荒诞不经。」见杨守敬、熊会贞疏，《水经注疏》，（江苏：古籍出版社，1989)，页3。}
\BookNote{8.}{《水经注》所引为《山海经·北山经》之文。}
\end{BookNotes}

% 原书第 144 页（PDF 第 157 页）
\BookPage{144}{157}
RHia碳空

楼南海量

薛海

新器定

\begin{BookQuote}
民尼蟒依阿林谷
\end{BookQuote}

静水AARG福祥察地注衡安

《水经注》河源图，采自《钦定河源纪略》土蕃、昆仑之方位，《新唐书》资料从《旧唐书》来，但比《旧唐书》详细；但两者之间对吐蕃距离中土的估计有些误差，因此胡渭质疑：按《旧唐书》吐蕃在长安西八千里，盖以其都城言之。今元鼎所见紫山，东距长安裁五千里，度其地望，当在都城三千里矣。吐蕃自言昆仑山在国中西南，而元鼎所指以为西南者，乃在其东北，与吐蕃所谓昆仑者，相去悬绝，是大可疑也。10

《新、旧唐书》对于昆仑山方位距离的判定出现矛盾。

对于昆仑为河源的记载，《史记》也有说法；〈大宛列传〉云：

于窴之西，水皆西流，注西海；其东则东流，注盐泽。盐泽潜行地下，其南则河源出焉。多玉石，河注中国。而楼兰、姑师邑有城郭，临盐泽。盐泽去长安可五千里。

\vfill
\begin{BookNotes}
\BookNote{9.}{张羽新主编，《中国西藏及甘青川滇藏区方志汇编》，（北京：学苑出版社，2003)，第51册，卷三，页34-35。}
\BookNote{10.}{参胡渭着、邹逸麟整理，《禹贡锥指》卷十三上，（上海：古籍出版社，1996版），页434。}
\end{BookNotes}

% 原书第 145 页（PDF 第 158 页）
\BookPage{145}{158}
\noindent 同〈传〉又曰：

\begin{BookQuote}
汉使穷河源，河源出于寘，其山多玉石，采来，天子案古图书，名河所出山曰
\end{BookQuote}

昆仑云。

司马迁于本〈传〉后的「赞」曰:

\begin{BookQuote}
禹本纪言「河出昆仑。昆仑其高二千五百余里，日月所相避隐为光明也。其上有醴泉、瑶池」。今自张骞使大夏之后也，穷河源，乌睹〈本纪〉所谓昆仑者乎？故言九州山川，《尚书》近之矣。至《禹本纪》、《山海经》所有怪物，
\end{BookQuote}

余不敢言之也。

司马迁以「疑则存疑」的态度探索河源，虽没强烈批判《山海经》提出的河源位置；但还是认为尽信古籍而不辨真伪不可取。王充推敲司马迁之意，认为太史公暗讽《山经》、《禹纪》为虚妄之言。《论衡·谈天篇》谓：

\begin{BookQuote}
太史公曰：「《禹本纪》言河出昆仓，…，余不敢言也。」夫弗敢言者，谓之虚也。昆仑之高，玉泉、华池，世所共闻，张骞亲行无其实。案《禹贡》九州山川怪奇之物、金玉之珍，莫不悉载，不言昆仑山上有玉泉、华池。案太史公之言，《山经》、《禹纪》，虚妄之言。
\end{BookQuote}

昆仑山河源的判断分持不同的看法，使得河源的问题更形复杂。张星烺认为或因历史上「昆仑」有两处，因此遂滋生这些误解，他说:「昆仑何在？我国学者自昔既有二说。《山海经》卷二〈西山经〉昆仑，毕沅注云在今甘肃肃州南八十里，又金城临羌县有昆仑祠，敦煌广至县有昆仑帐。此其一也。《史记·大宛列传〉『汉使…』；《尔雅·释地》云『河出昆仑』。…凡此记载皆误信黄河之源出于于阗，流入塔里木，至罗布泊尔，再流经地下至青海复出者。昆仑不在肃州而远在于阗，此第二说也。河源虽误而昆仑地位不误，数千年来，相传如此。萧州金城敦煌之昆仑谅因偶尔同名，故遂误会也。」11看来「河出昆仑」的历史认识，似乎需要重新审视考虑。

对于河源，《汉书·西域传》一开始就点出：

\begin{BookQuote}
河有两源：一出葱岭山，一出于阗。于阗在南山下，其河北流，与葱岭河合，
\end{BookQuote}

\vfill
\begin{BookNotes}
\BookNote{11.}{参张星烺，《中西交通史料汇编》，(北京:中华，1977)，第1册，页15。}
\end{BookNotes}

% 原书第 146 页（PDF 第 159 页）
\BookPage{146}{159}
东注蒲昌海。蒲昌海，一名盐泽者也，去玉门、阳关三百余里，广袤三百里。其水亭居，冬夏不增减，皆以为潜行地下，南出于积石，为中国河云。《汉书》的意见可以再议；但他所提出两源的说法，大抵为历来讨论河源最基本的观点。《元史·地理志》〈河源附录〉曰:

河源古无所见。《禹贡》导河，止自积石。汉使张骞持节，道西域，度玉门，见二水交流，发葱岭，趋于阗，汇盐泽，伏流千里，至积石而再出。唐薛元鼎使吐蕃，访河源，得之于闷磨黎山。然皆历岁月，涉艰难，而其所得不过如此。世之论河源者，又皆推本二家。其说怪迂，总其实，皆非本真。意者汉、唐之时，外夷未尽臣服，而道未尽通，故其所往，每迂回艰阻，不能直抵其处而究其极也。

《元志》明显承受《汉书》说法的影响。

\BookArt{p159-01.jpg}{0.891}{0.407}

\noindent 对于河源的描述，《尚书·禹贡》的「导河积石」也是历来集中讨论的焦点，

\vfill
\begin{BookNotes}
\BookNote{12.}{张羽新主编，《中国西藏及甘青川滇藏区方志汇编》，（北京：学苑出版社，2003)，第51册，卷三，页33。}
\end{BookNotes}

% 原图说　《汉书》河源图，采自《钦定河源纪略》

% 原书第 147 页（PDF 第 160 页）
\BookPage{147}{160}
\noindent 〈禹贡〉云：

\begin{BookQuote}
导河积石，至于龙门，南至于华阴，东至于底柱，又东至于孟津，东过洛汭，至于大伾；北过降水，至于大陆，又北播为九河，同为逆河，入于海。
\end{BookQuote}

「积石」，《汉书·地理志》「河关」下云：

\begin{BookQuote}
积石山在西南羌中，河水行塞外，东北入塞内。至章武入海，过郡十六，行
\end{BookQuote}

九千四百里。

《汉书·西域传》又云：

\begin{BookQuote}
盐泽水潜行地下，南出于积石，为中国河。
\end{BookQuote}

河源不在此处，众家皆谓禹在积石导河水流向龙门；因此，《书·大传》及〈正义》都以为此处为大禹施功之处也，如〈禹贡〉「雍州」所谓「浮于积石，至于龙门西河，会于渭汭。」而《汉志》所谓「积石」在河关县西南羌中；汉河关县，宣帝神爵二年（前60）置，十六国时期相继为割据政权控制，其中吐谷浑盘踞的时间最久。

但是，《通典》卷一四七〈州郡四〉云：

\begin{BookQuote}
按禹导河积石，…以此施功发迹，自积石山而东，即今西平郡龙支县界山是
\end{BookQuote}

也。如此便出现两处积石山的位置了。「龙支」为汉允吾县，属金城郡。13杜佑于此，似乎缺乏明辨。从《新、旧唐书》所载积石的位置，与杜佑所说的略有出入；《旧唐书·西戎传》「吐谷浑」载：

\begin{BookQuote}
(贞观九年）侯君集与江夏王道宗趣南路，登汉哭山，饮马乌海，获其名王梁屈忽，经涂二千余里空虚之地，盛夏降霜，多积雪，其地乏水草，将士啖冰，马皆食雪。又达于柏梁，北望积石山，观河源之所出焉。
\end{BookQuote}

\vfill
\begin{BookNotes}
\BookNote{13.}{参李吉甫撰、贺次君校，《元和郡县图志》卷三十九〈陇右道上〉，（北京：中华，1983），页989、页992。}
\end{BookNotes}

% 原书第 148 页（PDF 第 161 页）
\BookPage{148}{161}
\noindent 此事《新唐书·西域传》「吐谷浑」所载亦同于《旧唐书》。14因此，《元和郡县图志》对于积石山的叙述，或为可信。《元和志》卷三十九〈陇右道上〉，曰：

\begin{BookQuote}
积石山，一名唐述山，今名小积石山，在县西北七十里。按河出积石山，在西南羌中，注于蒲昌海，潜行地下，出于积石，为中国河，故今人目彼山为大积石，此山为小积石。
\end{BookQuote}

其后，同书本卷在「龙支县」下提及积石山，谓：

\begin{BookQuote}
积石山，在县西九十八里。南与河州枹罕县分界。
\end{BookQuote}

此即或为杜佑所谓龙支县辖下的另一处积石山耳。胡渭说：

\begin{BookQuote}
应劭云：析支东去河关千余里，河首积石，南枕析支。则县距此山亦千余里可知矣。自东晋之后，县为吐谷浑所据，遂以枹罕为华戎之界。故唐初太子贤注《后汉书》，误认龙支之积石为〈禹贡〉之积石。15
\end{BookQuote}

或如《元志》、〈正义〉所云，积石山有大、小积石之分，殆无疑义，〈正义〉云：

\begin{BookQuote}
河自盐泽潜行入吐谷浑界大积石山，又东北流至小积石山。
\end{BookQuote}

至此，〈禹贡〉所谓「导河积石」的位置隐隐若现。

或因经文「导河积石」的深刻印象，但河源又不是从积石山始；因此，长期以来对于河水的认识，存在一个错误印象，即以为河源从昆仑山起，流入渤泽，然后在积石山下又冒出。这是长期以来误以为黄河在河源的河段，经过「伏流重源」的现象所致。《山海经·西山经》曰：

\begin{BookQuote}
不周之山，…东望渤泽，河水所潜也，其源浑浑泡泡。
\end{BookQuote}

其后，《史记》对于河源「伏流重源」的描述，更使这个认识深植于一般的观念中。《史记·大宛列传》基本上是司马迁根据张骞出使西域的观察记录的，而张骞对于河源的认识，与实际状况有出入。〈大宛列传〉谓：

\vfill
\begin{BookNotes}
\BookNote{14.}{《新唐书·西域传》「吐谷浑」载：「君集、道宗登汉哭山，战乌海，获名王梁屈葱。…君集、道宗行空荒二千里，盛夏降霜，乏水草，士糜冰，马秣雪。阅月，次星宿川，达柏海上，望积石山，览观河源。」}
\BookNote{15.}{见胡渭着、邹逸麟整理，《禹贡锥指》卷十三上，(上海:古籍出版社，1996)，页410。}
\end{BookNotes}

% 原书第 149 页（PDF 第 162 页）
\BookPage{149}{162}
\begin{BookQuote}
于窴之西，则水皆西流，注西海；其东水东流，注盐泽。盐泽潜行地下，其南则河源出焉。多玉石，河注中国。而楼兰、姑师邑有城郭，临盐泽。盐泽去长安可五千里。匈奴右方居盐泽以东，至陇西长城，南接羌，隔汉道焉。
\end{BookQuote}

张骞两度被匈奴所拘，都在匈奴居盐泽以东至陇西长城一带；因此，他是否见到盐泽水文水域的分布状况，可疑。此或他在途中听闻。16

前引《汉书·西域传》提出河水两源的说法，更使河源「伏流重源」的说法，变成难以撼摇的历史认识。《水经·河水一》载「伏流重源」，谓：

\begin{BookQuote}
余考群书，咸言河出昆仑，重源潜发，沦于蒲昌，出于海水。
\end{BookQuote}

「伏流重源」的臆说影响非常深远。这是以推测的想像，将《山海经》、《尚书》以来对河源描述，作最方便的解释。17陶宗仪对此稍有质疑，《南村辍耕录》谓：

\begin{BookQuote}
史称河有两源，一出于于阗，一出葱岭。于阗之水北行，合葱岭河，注蒲类海(即蒲泽；下同)，不流，洑至临洮出焉。今洮水自南来，非蒲类明矣。询之土人，言于阗、葱岭水，其下流散之沙自碛。18
\end{BookQuote}

《辍耕录》引元代实地勘察而写成的《河源志》。不过，陶氏的质疑诘难，对「伏流重源」似乎影响不大。19直到清中、晚期，但仍有「伏流重源」的误解。清代祁韵

\vfill
\begin{BookNotes}
\BookNote{16.}{参《史记·大宛列传》。张骞去程系「与堂邑氏胡奴甘父俱出陇西。经匈奴，匈奴得之，传诣单于。单于留之。」而返程则「欲从羌中归，复为匈奴所得，留岁余。」此系岑仲勉的怀疑，合理；姑引之。参岑仲勉，《黄河变迁史》，（台北：里仁书局，民71翻印），页32-33。}
\BookNote{17.}{参郑德坤，〈层化的河水流域地名及其解释〉，收于氏着《中国历史地理论文集》，（台北：联经，民70)，页179-263。}
\BookNote{18.}{参陶宗仪，《南村辍耕录》卷二十二〈黄河源〉，（北京：中华，1997重刷），页268。丁谦在《汉书西域考证》中谓：「河水潜行，复出积石，古说相传已久，惟西人不信之，谓用实测法，测得罗布泊高于海平线二千六百尺，鄂陵泊高一万四千尺，至潜源重出之噶达素斋老更高至一万四千七百尺，水即能潜行千百里之远，岂能上涌千数百丈之高，此亦古今一大问题矣。」引自郑德坤，〈层化的河水流域地名及其解释〉，收于氏着《中国历史地理论文集》，（台北：联经，民70)，页194-195。}
\BookNote{19.}{另外，《通典》卷一七四〈州郡四·古雍州下·西平郡〉指出「潜流地下，…终是纰谬」。《舆地广纪》卷十六〈陕西秦凤路·积石军〉更直言：「河出昆仑，自古言者皆失其实，《禹本纪》、《山海经》、《水经》固以迂怪妄诞，而班固所载张骞穷河源事，亦为臆说。骞使大夏，见葱岭、于阗二河合流于蒲昌海，其水亭居，皆以为潜行地中，南出于积石为中国河，此臆度之，非实见蒲昌海与积石河通流之。}
\end{BookNotes}

% 原书第 150 页（PDF 第 163 页）
\BookPage{150}{163}
\noindent 士（1751-1815）是近代中国西北地理学研究的开创者，在其《西陲要略》中，说：「古之所谓河出昆仑墟，盖因天山、葱岭及穆苏尔冰岭一带，蒙古语呼之为昆都仑，急呼则为昆仑。有罗卜诺尔（即前之盐泽、蒲昌海；今称罗布泊）所汇之水，无非自天山、葱岭而来，特由罗布诺尔伏流地中，至噶达苏赤老始见，是为黄河之源，或以为罗卜诺尔为即星宿海，误矣。」20祁氏的错误见解流传在当时的知识界，也见于当代其他著作中。21

\BookArt{p163-01.jpg}{0.800}{0.448}

\vfill
\begin{BookNotes}
\BookNote{20.}{见祁韵士，《西陲要略》卷一〈南北两路山水总叙〉，收于氏着《万里行程纪》（外五种），（山西：人民出版社，1992），页134。}
\BookNote{21.}{清光绪年间，陶保廉在著名的《辛卯侍行纪》中，仍有对河源判断错误的说法，陶氏说：「河有重源，均出昆仑，稽古证今，一一吻合。」足见「重源说」的传统，对当时影响仍颇具影响。引自田尚，〈黄河河源探讨〉，收于祁明荣编《黄河源头考察文集》，（青海：人民出版社，1982），页83。}
\BookNote{22.}{上海古籍出版社编，《宋元笔记小说大观》，（上海：上海古籍出版社，2001)，第6册，卷二十二，页6412。}
\end{BookNotes}

% 原图说　《南村辍耕录》所附〈黄河源〉图

% 原书第 151 页（PDF 第 164 页）
\BookPage{151}{164}
\BookSubheading{三、河源探勘：历史上探源的行动}

贞观九年（635）侯君集等人抵达河源，是文献上记载最早到此活动的汉人，《旧唐书·西戎传》「吐谷浑」载：

\begin{BookQuote}
(贞观九年；635)侯君集与江夏王道宗趣南路，登汉哭山，饮马乌海，获其名王梁屈忽，经涂二千余里空虚之地，盛夏降霜，多积雪，其地乏水草，将士啖冰，马皆食雪。又达于柏梁，北望积石山，观河源之所出焉。
\end{BookQuote}

《旧唐书·侯君集传》所载略同，惟增加了「转战过星宿川，至于赤海…北望积石山，观河源所出焉。」。《新唐书·西域传（上）》「吐谷浑」则载王道宗语，谓：「柏海近河源，古未有至者。」说明在此之前罕有人迹。柏海即赤海，亦称柏梁，即今之札陵湖。这是文献中较为明确记载汉人登河源的资料，也是文献中黄河源区出现湖泊的最早记载。另在《旧唐书·吐蕃传（上）》记吐蕃赞普弃宗弄赞求亲，载：

\begin{BookQuote}
贞观十五年(641），太宗以文成公主妻之，令礼部尚书、江夏郡王道宗主婚，持节送公主于吐蕃。弄赞率其部兵次柏海，亲迎于河源。
\end{BookQuote}

这些是关于河源实际观察的较早资料。

唐代第二次抵河源的记载，为穆宗长庆二年（821）刘元鼎出使吐蕃，《旧唐书·吐蕃传（下）》谓：

\begin{BookQuote}
元鼎往来蕃中，…往来，渡黄河上流，在洪济桥西南二千余里，其水极为浅狭，春可揭涉，秋夏则以船渡。其南三百余里有三山，山形如鏊，河源在其间，水甚清冷，流经诸水，色遂赤，续为诸水所注，渐既黄浊。又其源西去蕃之列馆约四驿，每驿约二百余里。东北去莫贺延碛尾，阔五十里，向南渐狭小，北自沙州之西。乃南入吐浑国，至此转微，故号碛尾。计其地理，当剑南
\end{BookQuote}

之直西。

新、旧《唐书》所载元鼎使蕃事略同，《新唐书》稍详；前引，故不赘。但从新、旧《唐书》所载，河源为入西藏吐蕃「中高而四下」的三山。康熙四十三（1704)年，舒兰奉旨探勘河源，返程后呈〈河源记〉一篇，曰：


% 原书第 152 页（PDF 第 165 页）
\BookPage{152}{165}
\BookArt{p165-01.jpg}{0.990}{0.426}

（六月）初九至星宿海，蒙古名鄂敦塔拉择山，至高者登而视之，观星宿海之源，小泉万亿不可胜数；周围群山蒙古名为库尔棍，即昆仑也。南有高山名古而班吐而哈，西有高山名巴尔布哈，北有高山名阿克塔因凄奇，三山之泉流出三支河，即古儿班索而嘛也。三河往东顺流入于扎陵，自扎陵一支流入鄂陵，自鄂陵流出之河乃黄河也。24

刘元鼎的观察与清人的实地考察吻合；刘元鼎的描述是经过实地勘察后，明确指出河源的最早记录，代表唐代探索河源的成就。

元代是历史上最早明确订出探勘河源而付诸行动的朝代，《元史·地理志》〈河源附录〉载：

至元十七（1280）年，命都实为招讨使，佩金虎符，往求河源。都实既受命，

\vfill
\begin{BookNotes}
\BookNote{23.}{张羽新主编，《中国西藏及甘青川滇藏区方志汇编》，（北京：学苑出版社，2003)，第51册，卷三，页36。}
\BookNote{24.}{参舒兰，〈河源记〉，收于《小方壶斋舆地丛钞》第四帙〈河源记〉一。}
\end{BookNotes}

% 原图说　唐刘元鼎所见河源图

% 原书第 153 页（PDF 第 166 页）
\BookPage{153}{166}
\begin{BookQuote}
是岁至河州。州之东六十里，有宁河驿。驿西南六十里，有山曰杀马关，林麓穹隘，举足浸高，行一日至巅。西去愈高，四阅月，始抵河源。是冬还报，并图其城传位置以闻。其后翰林学士潘昂霄从都实之弟阔阔出得其说，撰为《河
\end{BookQuote}

源志》。

都实的探源，稍后由潘昂霄整理出《河源志》一书；后来朱思本据此书梵文本，翻译成汉文，而成详、略二本。其中朱思本的汉文翻译本较详，而《元史·地理志》〈河源附录〉即将两种合为一编的足本。《河源志》的刊刻，对于黄河探源行动具有某些意义。柯九思谓：

\begin{BookQuote}
河源有志，自本朝始，前乎此，曷为未有。志河源者，道路辽阻，所传闻异辞，莫能究河之源也。
\end{BookQuote}

在〈河源附录〉载朱思本的「河源」，谓：

\begin{BookQuote}
河源在中州西南，直四川马湖蛮部之正西三千余里，云南丽江宣抚司之西北一千五百余里，帝师撒思加地之西南二千余里。水从地涌出如井。其井百余，东北流百余里，汇为大泽，曰火敦脑儿。
\end{BookQuote}

从火敦脑儿西南方向来注的河川，判断即今之喀喇渠；都实知道流百余里来注的喀喇渠，却仍将河源放在星宿，这是他的误判。「火敦脑儿」即〈河源志〉中所说的「星宿」，曰：

\begin{BookQuote}
按河源在土蕃朵甘思西鄙，有泉百余泓，沮洳散涣，弗可逼视，方可七八十里，履高山下瞰，灿若列星，以故名火敦脑儿。火敦，译言星宿也。群流奔辏，近五七里，汇二巨泽，名阿刺脑儿。
\end{BookQuote}

阿剌脑儿二巨泽，即今之扎陵湖与鄂陵湖；这是河源地区二大泽见于文献最早的记载。另外，都实探源最大的意义，在于厘清河源一些错误的传统观念，如「重源」的讹传。

目前现存最早的河源图，为元人陶宗仪《南村辍耕录》卷二十二〈黄河源〉所附的「黄河源」图；图中虽未标明是否与都实的探源相关，但所绘地名及河源附近的地理分布与《河源志》所述吻合。《辍耕录》在河源部分，仅绘至星宿海与《河源


% 原书第 154 页（PDF 第 167 页）
\BookPage{154}{167}
\noindent 志》同。

\BookArt{p167-01.jpg}{0.855}{0.504}

洪武十五年（1382），僧宗泐奉使西藏归来，经过河源地区，以〈河源诗〉自记，诗前有序，曰：

\begin{BookQuote}
河源出自抹必力赤巴山，番人呼黄河为抹处，牦牛河为必力处，赤巴者分界也。其山西南所出之水，则流入牦牛河，东北之水是为河源。…其源东抵昆仑可七八百里，所涉处尚三百余里，下与昆仑之水合流。26
\end{BookQuote}

宗泐的「河源诗序」是研究河源史一则重要的材料；对于河源距离的远近，可与唐

\vfill
\begin{BookNotes}
\BookNote{25.}{张羽新主编，《中国西藏及甘青川滇藏区方志汇编》，（北京：学苑出版社，2003)，第51册，卷三，页37-38。}
\BookNote{26.}{参《河源纪略》卷二十四〈辨讹五〉，宗泐〈河源诗自记〉。}
\end{BookNotes}

% 原图说　元使所穷河源图

% 原书第 155 页（PDF 第 168 页）
\BookPage{155}{168}
\noindent 代刘元鼎观察河源的资料互证。黄盛璋推测宗泐由藏归来的路线与元鼎一样。27这条入藏的通道，始于唐代所辟，一直到清代，基本上都探行此一路线；其路线自鄯州（西宁），经赤岭（日月山），取道黄河上源，出扎陵、鄂陵两湖间，沿喀喇渠，越黄河、长江分水岭而渡厘牛河，由此经当拉岭入拉萨。28抹必力赤巴山即今巴颜喀喇山，它也是黄河与厘牛河的分水岭；厘牛河即今通天河，为长江上源。抹处的「处」，即藏语「河」的意思，而「抹处」即藏语所称的「玛曲」。宗泐的〈河源诗〉代表了明代对河源的认识。

清代对于河源的探勘较为积极，前后数次进行探源行动；29仅康熙年间（1662-1723）就先后两次派员勘察河源。康熙四十三年（1704），拉锡、舒兰等奉命前去河源考察。返京后，撰写〈河源记〉及绘制〈星宿河源图〉。30

\begin{BookQuote}
<河源记〉前引，兹不录。文中除了描述星宿地区的地形状况，同时也记载河
\end{BookQuote}

源系出三河；〈河源记〉记载黄河上源分为三支，皆东流入扎陵湖；其中在〈星宿河源图〉中渣凌诺尔（扎陵湖）西绘有三支河，中间一条绘得特别长，河上注记：「黄河源三河名古儿班索尔嘛」；图上注记与〈河源记〉的文字相同。〈河源记〉与〈星

\vfill
\begin{BookNotes}
\BookNote{27.}{参黄盛璋，〈再论黄河河源问题〉，收于氏着《历史地理论丛》，（北京：人民出版社，1982)，页321-331。黄盛璋说:「宗泐渡河之处，距河源尚有三百余里；元鼎所渡河，河源也是在其南三百余里，情形完全相同。」}
\BookNote{28.}{参《新唐书》卷四十地理志〉「鄯州」条载抬公主入藏路线，「鄯州，西平郡…；又西二十里至赤岭，其西吐蕃，有开元中分界碑。自振武经尉迟川、苦拔海、王孝杰米栅九十里至莫离驿。又经公主佛堂、大非川二百八十里至那录驿，吐浑界也。…。又渡西月河，二百一十里至多弥国西界。又经声牛河度藤桥，百里至列驿。…乃度大月河罗桥，经潭池、鱼池，五百三十里至悉诺罗驿。…唐使入蕃，公主每使人迎劳于此。…又六十里至突录济驿，唐使至，赞普每遣使慰劳于此。…逻些在东南，距农歌二百里，唐使至，吐蕃宰相每遣使迎候于此。」〈地理志〉清楚地指出唐代入藏的路线。}
\BookNote{29.}{王尧认为清初的探源行动，为清代的河源知识奠下基础，后来出现大量作品，如《大清一统志》(1743)、《水道提纲》(1761)、《河源纪略》(1782)、《星宿河源图》(1708)、《一统志青海图》（1717)、《皇舆全览图》（1718)、《乾隆内府舆图》（又称《乾隆十三排图》，1760），皆与清初的探源有关。参王尧，〈黄河源上两大湖—扎陵、鄂陵位置考实〉，收于祁明荣编《黄河源头考察文集》，（青海：人民出版社，1982），页136。}
\BookNote{30.}{〈星宿河源图〉长年藏于北京故宫，外界莫由得知。后来才在故宫明清档案馆中发现绢本彩绘。其图注记汉文用黄纸，满文用白纸贴，后来又加红色注记。从图的样式、装潢、写绘，被认为是进呈御览。关于〈星宿河源图〉与〈河源纪〉两者关系，可参黄盛璋，〈黄河上源的历史地理问题与测绘的地图新考〉，收于氏着《历史地理论丛》，(北京：人民出版社，1982)，页332-362。}
\end{BookNotes}

% 原书第 156 页（PDF 第 169 页）
\BookPage{156}{169}
\noindent 宿河源图〉较以往绘制河源支河分布的地图详细甚多，在探讨河源的汉文古籍中，最早出现「扎陵湖」、「鄂陵湖」名称及指出具体位置，以拉锡、舒兰等奉命前去探源后的记录。《康熙朝东华录》康熙四十三（1704）年九月丁卯，载：

\begin{BookQuote}
上遣侍卫拉锡等探视河源，谕之曰：黄河之源，名古尔班所罗谟，其实发源之处，从来无人到过。尔等务须直穷其源明白，视察其河流至何处入雪山边。…拉锡等回奏：…臣等起程，前行六月七日至星宿海之东，有泽名鄂陵，周围一百余里；初八日，至鄂陵西有泽名扎陵，周围三百余里。…星宿海之源，小泉万亿不可胜数。.31
\end{BookQuote}

\BookArt{p169-01.jpg}{0.990}{0.538}

\vfill
\begin{BookNotes}
\BookNote{31.}{参《康熙朝东华录》卷十五，康熙四十三年九月丁卯，（台北：大东书局，出版年代不详），页575。}
\BookNote{32.}{曹婉如等编着，《中国古代地图集·清代》，（北京：文物出版社，1997)，页2。图藏中国第一历史档案馆。}
\end{BookNotes}

% 原图说　〈星宿河源图〉

% 原书第 157 页（PDF 第 170 页）
\BookPage{157}{170}
将〈河源记〉所载，对照乾隆二十六年（1761）齐召南所编的《水道提纲》记载的河源，当可更明确了解清代对河源的认识越来越清晰。《水道提纲》卷五〈黄河〉载：

\begin{BookQuote}
黄河源出星宿海，巴颜喀喇山之东麓，二泉流数里，合而东南，名为阿尔坦河。…东流数十里，折东北流百里，至鄂端塔拉即古星宿海，《元史》所谓鄂端诺尔也，自河源至此已三百里。星宿海于群山围绕中，平地有泉千百泓，并涌望若列星。阿尔坦河自西南来…，北有巴尔哈布山，西南流出一水；南有哈喇荅尔罕山北流出之水来会为一道。33
\end{BookQuote}

齐召南未亲临河源探勘，但《水道提纲》所描述的河源河道，是清代最详细的，也是后来讨论河源重要的材料。

康熙五十六年（1717），派喇嘛楚儿沁藏布、兰木占巴、理藩院主事胜住等前往青海、西藏绘制舆图，对于河源地区务求「一山一水，悉入图志」。这一行人实地测绘的地图，交给了传教士雷孝思（Jean-BaptisteRegis,1663-1728）等人编绘〈皇舆全览图〉作参考。《水道提纲》的〈原序〉推崇「圣祖（康熙）御制舆图非前史志所能稍及」，齐召南受〈全览图〉影响绝非不可能。

乾隆四十七年（1782），帝命阿桂之子阿弥达「恭祭河源」，此事记于《河源纪略》卷首，文曰：

\begin{BookQuote}
今年春间豫省青龙岗漫口合龙未就，遣大学士阿桂之子干清门侍卫阿弥达前往青海，务穷河源，告祭河神。事竣复命，并据按定南针绘图具说呈览。具奏星宿海西南，有一河名阿勒坦郭勒，蒙古语阿勒坦即黄金，郭勒即河也。此河实系黄河上源，其水色黄迥旋三百余里，穿入星宿海。
\end{BookQuote}

阿弥达访河源，系为了告祭河神；探源祇是此行的附带。因此，他的穷河源并未确实。如在他所进呈的河源图〉中，错置了扎陵湖为「阿拉克淖尔」、鄂陵湖为「策克淖尔」。其实「阿拉克淖尔」在康熙时期的〈星宿河源图〉中已出现，在扎陵湖之北。而扎陵、鄂陵两湖名称沿用至今。另外，还继续沿用「黄河重源」等错误的说

\vfill
\begin{BookNotes}
\BookNote{33.}{参齐召南，《水道提纲》卷五〈黄河〉，文渊阁《四库全书·史部》十一，（台北:台湾商务印书馆影印，民75)。}
\end{BookNotes}

% 原书第 158 页（PDF 第 171 页）
\BookPage{158}{171}
\noindent 法。34

不过，阿弥达访河源后，乾隆还命四库馆臣将他们的考察记录汇编而成《河源纪略》一书，显见乾隆对阿弥达访河源的重视。

\BookArt{p171-01.jpg}{0.990}{0.534}

\vfill
\begin{BookNotes}
\BookNote{34.}{参纽仲勋，〈黄河河源考察和认识的历史研究〉，载于《中国历史地理论丛》，1988年，第4辑。}
\BookNote{35.}{美国国会图书馆〈河源图〉，登录号:G7822.Y4A51795.A5Vault°}
\end{BookNotes}

% 原图说　满、汉文对照的〈星宿河源图〉

% 原书第 159 页（PDF 第 172 页）
\BookPage{159}{172}
\BookSubheading{四、小结}

从院藏数幅黄河图对于河源地区的描述，引发个人对于河源探讨的兴味；文中并非寻找新的事证来重新界定河源，全文旨在对于黄河源头的诸多历史认识，作进一步的了解。

从先秦以来对于河源的认识，经过漫长的摸索探勘，最后方如剥笋般层层揭开河源的真相。历史上最初对于河源的解释，是透过经书古籍的记载，作一些臆断式的推测；其后，陆续有人登上河源地区作一些勘察报告。或因探勘地点的不同，因此对于河源地区的描述也不尽相同。不过，从历史上认识河源的过程，我们大致可以界分成三个阶段：一、先秦以来，从传闻或文献上来讨论河源，其实当时并没有人登上河源地区。二、两汉到魏晋、隋唐间，若干出使的使臣或征战的军士，并非为了探源而登上河源附近地区，他们将亲历的状况作一些描述；不过，此时还是在文献上寻求材料来证明所观察到的河源。三、从元迄明、清，已目标明确的为了探源而登入河源地区，作接近河源真实状况的描述了。

其实河源的认识与确定，一直到二十世纪中叶，仍出现一些争议；如对于河源地区的各支流，如何判断主要的河源。以「源远流长」的河源，来寻求河源，这是学术上的共识；因此，在河源地区一直存在黄河源支流「约古宗列曲」与「卡日曲」孰长孰短的争议。不过，这不是本文处理的问题，搁置不论。透过现代科技与卫星航照的探索，或许能完全解决河源的问题；但本文主要的目的是在于讨论一项「历史认识」的过程。透过这个过程，我们看到了历史上各个时期对于河源一步步的摸索。


% 原书第 160 页（PDF 第 173 页）
\BookPage{160}{173}
